% RESEARCH_DRAFT_SHA256: a01967ba1ad398f3a21adf5a50911f6ff3fb740bb4ec97f27c2c39f9436522c5
% Layout and source-reference wrappers expanded for the web renderer.
\documentclass[11pt]{article}
\usepackage[margin=25mm]{geometry}
\usepackage{fontspec,unicode-math}
\setmainfont{Latin Modern Roman}
\setmathfont{Latin Modern Math}
\usepackage{amsmath,amssymb,mathtools,xurl,hyperref}
\hypersetup{colorlinks=true,urlcolor=blue,linkcolor=blue}
\setcounter{secnumdepth}{0}
\setlength{\parindent}{0pt}
\setlength{\parskip}{5pt}
\setlength{\emergencystretch}{3em}
\newcommand{\sref}[2]{\hyperlink{source:#1:#2}{[S#2]}}
\newcommand{\cataloguescope}[1]{\paragraph{Source record.}#1}
\newcommand{\floor}[1]{\left\lfloor#1\right\rfloor}
\begin{document}
\section*{Research-attempt provenance}
This standalone research attempt retains the complete original definition body below,
including its historical source and preparation notes. Those notes describe the
earlier catalogue preparation; the new attempt and its verification are identified
separately after that body. No original ID, catalogue status, or source statement is changed.

\section{Expansion-Point Selection in Time-Domain Model Reduction}\label{statement}

\subsection{Definitions and mathematical statement}

\paragraph{Shared notation.}Standard number systems are $\mathbb N,\mathbb Z,\mathbb Q,\mathbb R,\mathbb C$. Specialized notation and index ranges follow the cited source and the definitions below.


The impulse response is the output generated by an impulse input. Rational Krylov reduction uses shifted resolvents around expansion points; a Laguerre basis represents time responses using functions with a chosen scale $\alpha$.
\paragraph{Record scope and qualifications.}
The catalogue title is retained as metadata. Named certificates, contracts, bounds or frameworks in that title are not asserted as results of the original source extract.
The optimization norm, time horizon, admissible parameters, and input-system assumptions are not stated.

Source-record limitation: The expansion-parameter selection target has no specified error norm, time horizon, or optimization criterion.
\paragraph{Statement recorded in the supplied source.}
In rational Krylov or Laguerre-based time-domain model reduction, how should the expansion point $s_0=\alpha$, equivalently the Laguerre parameter $\alpha$, be chosen to approximate the impulse response accurately? \href{https://doi.org/10.4171/owr/2009/11}{[S2]}
\paragraph{Source status.}
The supplied catalogue labels this source record partially solved; that metadata was checked against the supplied snapshot on 2026-10-11. The label is preserved for provenance and is not an independent verification of present open status or mathematical correctness. The catalogue literature assessment is dated 2026-08-21.
\subsection{Short English statement}
Choose the expansion or Laguerre parameter to approximate impulse responses accurately.
\subsection{Sources}
\begin{itemize}
\item[S1] Ulam AI and the UnsolvedMath Contributors, UnsolvedMath: A Curated Collection of Open Mathematics Problems, catalogue entry. Catalogue curation and metadata: CC BY 4.0. The source extract is retained; definitions, English summaries and qualifications are editorial additions.
\url{https://huggingface.co/datasets/ulamai/UnsolvedMath}.
\item[S2] Control Theory: On the Way to New Application Fields (2009), Oberwolfach Reports, DOI 10.4171/owr/2009/11, pp. 593--680
\url{https://doi.org/10.4171/owr/2009/11}.
\item[S3] Grune--Worthmann 2008
\url{https://skoge.folk.ntnu.no/prost/proceedings/ifac2008/data/papers/2302.pdf}.
\item[S4] shift-selection study
\url{https://doi.org/10.1007/s11044-022-09872-7}.
\item Source-specific Oberwolfach Reports licences, including CC BY-SA 4.0 where applicable, are retained; see THIRD\_PARTY\_NOTICES.md.
\end{itemize}
\paragraph{Preparation and review.}
The original source and UnsolvedMath contributors retain attribution for the source extract. Codex (GPT-6; exact serving revision not exposed) assisted with definitions, English summaries, scope notes and formatting. No new human mathematical authorship or independently verified solution is asserted. Catalogue identity checked against the complete supplied registry and current repository. Distinct upstream records, including repeated questions and broad or damaged extracts, are retained for complete coverage of the source CS category. A repeated record is not a new independent conjecture. Complete mathematical-equivalence and current-literature review has not been performed. Preparation and final review follow contribution prompts C and D, with source gaps disclosed. Structural and rendering checks do not establish a new theorem or open status.
\label{end:statement}

\section{Literature and context check}
The primary report on printed pages 637--638 already proposes a weighted Laguerre-coefficient objective and an analytic scale involving two impulse-response moments, followed by a reduced-model iteration. Those source claims are context, not new results or a convergence theorem rederived here. The following instead makes a specific auxiliary objective completely explicit: best one-dimensional $L^2(0,\infty)$ projection onto the first normalized exponential basis function. This tests scale selection and the limits of a universal parameter.
\section{Research attempt: Exact scale selection for a one-term exponential approximation}

\paragraph{One decaying exponential}
For $\alpha>0$ set $\phi_\alpha(t)=\sqrt{2\alpha}e^{-\alpha t}$. Direct integration gives $\lVert\phi_\alpha\rVert_2=1$. For any real $h\in L^2$, expansion of the square shows that $c\phi_\alpha$ uniquely minimizes $\lVert h-c\phi_\alpha\rVert_2^2$ at $c=\langle h,\phi_\alpha\rangle$, with error $\lVert h\rVert_2^2-c^2$.

For $h(t)=e^{-\beta t}$, $\beta>0$, this gives
\[
 c=\frac{\sqrt{2\alpha}}{\alpha+\beta},\qquad
 E(\alpha)=\frac1{2\beta}-\frac{2\alpha}{(\alpha+\beta)^2}
   =\frac{(\alpha-\beta)^2}{2\beta(\alpha+\beta)^2}.
\]
Thus the unique zero-error scale is $\alpha=\beta$, and the relative squared error is $((\alpha-\beta)/(\alpha+\beta))^2$. For any fixed chosen scale $\alpha$, this relative error approaches one as $\beta/\alpha$ tends to zero or infinity. Hence no system-independent positive scale gives a uniform relative squared error strictly below one even on these single-pole stable systems.

\paragraph{An exact finite calculation for a two-scale positive mixture}
Take $h(t)=e^{-\beta t}+w e^{-\gamma t}$, where $0<\beta<\gamma$ and $w>0$. Define
\[
 F(\alpha)=\frac1{\alpha+\beta}+\frac{w}{\alpha+\gamma},\quad
 \lVert h\rVert_2^2=\frac1{2\beta}+\frac{2w}{\beta+\gamma}+\frac{w^2}{2\gamma}.
\]
The projected energy is $G(\alpha)=2\alpha F(\alpha)^2$, so minimizing error is maximizing $G$. It tends to zero at both positive-domain ends, is positive within, and has derivative
\[
 G'(\alpha)=2F(\alpha)\left(
 \frac{\beta-\alpha}{(\alpha+\beta)^2}
 +w\frac{\gamma-\alpha}{(\alpha+\gamma)^2}\right).
\]
This derivative is positive for $\alpha\leq\beta$ and negative for $\alpha\geq\gamma$, with the strict endpoint signs because $w>0$ and $\beta<\gamma$. Therefore every global maximizing scale lies in $(\beta,\gamma)$ and is a root of the explicit cubic
\[
 (\beta-\alpha)(\alpha+\gamma)^2
 +w(\gamma-\alpha)(\alpha+\beta)^2=0.
\]
Its leading coefficient is $-(1+w)$, so it is not the zero polynomial. Evaluating $G$ at its finitely many real roots in the interval and choosing the greatest value gives an exact characterization of the optimum; for rational inputs this can be implemented with real algebraic root isolation and comparison. Alternatively, the formulas alone specify a finite list of candidates without making an implementation-complexity claim. Uniqueness of this optimizer is not assumed or proved.

\paragraph{The gap to the source's reduction problem}
This objective optimizes one orthogonal exponential approximation to the impulse response. Moment matching in a rational Krylov reduced state-space realization has additional algebraic constraints, and its finite-order approximation error need not equal $E(\alpha)$. Nor does reducing the impulse-response $L^2$ error alone prove a chosen norm bound for every input response. For example, scaling the input scales any nonzero linear response error. The two-scale stationary equation also depends on the actual mixture, so an iteration that estimates its parameters from a reduced model needs a separate convergence argument. The report's observed rapid iteration is not a proof of that argument.

\section{Verification and adversarial audit}
The basis norm, projection coefficient, error identity, derivative and cubic are derived explicitly. Equal scales reduce to the one-exponential case after an amplitude factor; $w=0$ does likewise. The strict positive-mixture assumption is used for the derivative signs and is not silently extended to signed modal residues. The limits $\alpha\downarrow0$ and $\alpha\uparrow\infty$ exclude boundary optima in the mixture calculation.
\paragraph{Audit corrections.} Acknowledged the parameter proposal already present in the primary source; distinguished its objective from the explicitly chosen auxiliary projection error.
\paragraph{External computation.} None; symbolic arguments and boundary checks only.
\section{Scope and remaining gaps}
A precise scale-selection problem is solved for one-term exponential projection, including a finite candidate rule for a positive two-mode family. General rational Krylov/Laguerre reduction, the original unspecified accuracy criterion, and convergence or optimality of the source's iterative parameter rule remain unresolved here.
\paragraph{Integration review.} Codex (GPT-6; exact serving revision not exposed) reviewed the written arguments and expanded layout and source-reference wrappers. Repository metadata and the production web renderer were checked using Python 3.13 and JavaScript with MathJax 3.2.2. This model-assisted review is not independent human proof verification; no PDF compilation or formal proof checking is asserted.

\section{Completion Estimate}
\noindent\textbf{COMPLETION: 0\%}
The calculations below the preserved source statement investigate a concrete route,
but establish no substantiated new progress toward the entire original target.
Known special cases and failed sufficient conditions are not credited as a solution.
\section{References}
The original source references and catalogue attribution are preserved in the
Sources section above. The following URLs were actually consulted for this attempt;
retention of an additional original reference does not claim fresh consultation of it.
\begin{itemize}
\item \url{https://ems.press/content/serial-article-files/46211} (consulted 2026-10-11). 
\end{itemize}
\section{Model and contributor disclosure}
The parent Codex GPT-6 model has reviewed the written mathematical arguments at their declared scopes. This is model-assisted review, not independent human verification.
OpenAI gpt-6-astra, requested reasoning configuration ultra, generated this research
attempt on 2026-10-11. The exact serving revision is not exposed. Codex GPT-6
performed orchestration and repository-format review; its exact serving revision
is also not exposed. The mathematical route, boundary checks and final adversarial
self-audit followed contribution prompts A then B in
\texttt{docs/contribution\_prompts.md}, after reading \texttt{CONTRIBUTING.md}.
Python 3.14.2 assembled the standalone documents; any mathematical computation is
specified in Verification. No theorem prover was used. No human mathematical
authorship, substantial human proof correction, or independent human proof checking
is asserted. Model-assisted review, source retrieval, and structural checks do not
establish correctness or novelty. Full private reasoning transcripts are unavailable
as reproducibility artifacts; the complete arguments, inputs and outcomes actually
used are stated here or in the companion batch audit.
\section{Final status}
Unresolved as a new research contribution. This is a concrete route investigation,
not a claim of new complete resolution or a current-status certification. A known
answer reproduced for context does not become a new contribution. The catalogue's
source-status label is preserved; this attempt status does not assert that a
source-labelled solved question is open.
\end{document}
